%================================================================
\chapter{Background}
\label{ch:background}
%================================================================

\section{Water Distribution Networks}
\label{sec:wdn_background}

\subsection{Hydraulic Fundamentals}

A Water Distribution Network (WDN) transports water from sources (reservoirs, tanks)
through a network of pipes to consumers at junctions. The network is
fully described by two sets of equations.

The \emph{continuity equation} enforces mass conservation at each
junction~$i$:
\begin{equation}
  \sum_{j \in \mathcal{N}(i)} q_{ij} = d_i,
  \label{eq:continuity}
\end{equation}
where $q_{ij}$ is the signed flow in the pipe connecting $i$ and $j$
(positive towards $i$) and $d_i$ is the demand at junction $i$. In
matrix form this reads $\mathbf{B}\mathbf{q} = \mathbf{d}$, where
$\mathbf{B} \in \{-1, 0, +1\}^{N \times M}$ is the node-edge
incidence matrix.

The \emph{head-loss equation} relates the pressure difference across
a pipe to the flow through it. For turbulent flow in a pipe of length
$L_{ij}$, diameter $D_{ij}$, and roughness $C_{ij}$, the
Hazen-Williams formula gives:
\begin{equation}
  p_i - p_j = r_{ij} \, |q_{ij}|^{1.852} \, \mathrm{sgn}(q_{ij}),
  \label{eq:headloss}
\end{equation}
where $r_{ij} = 10.67 L_{ij} / (C_{ij}^{1.852} D_{ij}^{4.87})$.
Solving equations~(\ref{eq:continuity}) and~(\ref{eq:headloss})
simultaneously for all junctions and pipes given boundary conditions
(reservoir heads, demand patterns) is the hydraulic simulation
problem. The EPANET solver~\cite{rossman2000epanet} is the
standard reference implementation; this thesis uses the Python
interface provided by WNTR~\cite{klise2017wntr}.

\subsection{Sensor Placement and Observability}

Practical networks instrument only a fraction of their junctions
and pipes. If the set of observed nodes is $\mathcal{V}_o \subset
\mathcal{V}$, the remaining junctions are \emph{unobserved}. From
the perspective of any learning model, the pressure readings at
unobserved nodes are missing values that must be predicted from
(a) the observed readings and (b) the graph structure.

This is a non-trivial interpolation problem. The pressure field is
smooth along pipes of similar roughness but can change abruptly at
high-demand nodes or at pressure-reducing valves. Classical
approaches try to recover the full pressure vector by state
estimation; machine learning approaches learn the mapping implicitly
from simulated data.

\subsection{EPANET Benchmark Networks}

The three networks studied in this thesis are standard EPANET
benchmarks widely used in the water systems literature.

\textbf{Net1}~\cite{rossman2000epanet} is a small reference network
with 11 junctions, 13 pipes, 1 reservoir, and 1 tank. Its small size
makes it a useful sanity-check but means it saturates quickly: even
a simple model achieves near-perfect reconstruction, leaving little
room for adversarial fine-tuning to improve matters.

\textbf{Net3}~\cite{rossman2000epanet} has 97 junctions, 117 pipes,
3 reservoirs, and 3 tanks. It includes two pumps and several
looped sections, so its flow field is more complex than Net1's.
Net3 has been used widely in the water-systems literature, including
GNN work on pressure reconstruction and state estimation; here it
serves as a mid-sized benchmark that sits between the small Net1 and
the large Modena network, which lets us see how the self-play results
scale with graph size.

\textbf{Modena}~\cite{bragalli2012modena} is a real-world network
from Modena, Italy, with 272 junctions and 317 pipes. It has no
pumps in the publicly distributed version, but its large, irregular
topology and heterogeneous demand patterns make it a challenging
test case.

Table~\ref{tab:networks} summarises the three networks.

\begin{table}[htbp]
  \centering
  \caption{Summary of the three benchmark networks.}
  \label{tab:networks}
  \begin{tabular}{lrrrl}
    \toprule
    Network & Junctions & Pipes & Source \\
    \midrule
    Net1    & 11  & 13  & EPANET reference~\cite{rossman2000epanet} \\
    Net3    & 97  & 117 & EPANET reference~\cite{rossman2000epanet} \\
    Modena  & 272 & 317 & Bragalli et al.~\cite{bragalli2012modena} \\
    \bottomrule
  \end{tabular}
\end{table}

\section{Cyber-Physical Attacks on Water Networks}
\label{sec:attack_background}

\subsection{Threat Landscape}

WDN SCADA systems face three categories of
digital threat: denial-of-service attacks that disrupt communication
links, integrity attacks that inject false data into the sensor
stream, and replay attacks that replay recorded legitimate data to
hide ongoing physical manipulation. Of these, integrity and replay
attacks are the hardest to detect because they do not disrupt
connectivity: from the network's perspective, the sensors appear to
be working normally.

The SCADA security community distinguishes \emph{targeted}
attacks, where an adversary selects sensors whose falsification
will have the greatest operational impact, from \emph{stealthy}
attacks, where the adversary's goal is to remain undetected for as
long as possible while still achieving an objective. These goals are
in tension: a large perturbation is easier to achieve physically but
easier to detect; a small drift is harder to detect but takes longer
to cause damage.

\subsection{False Data Injection Attacks}

A \emph{false data injection} attack replaces the true reading
$p_i^*$ at sensor $i$ with a corrupted value
$\tilde{p}_i = p_i^* + \delta_i$, where $\delta_i$ is the attacker's
perturbation~\cite{liu2011false}. If the detector relies on a linear
state-estimation residual $\mathbf{r} = \mathbf{z} - \mathbf{H}\hat{\mathbf{x}}$,
the attacker can construct $\boldsymbol{\delta} = \mathbf{H}\mathbf{c}$
for any vector $\mathbf{c}$, making the residual unchanged. This
attack is undetectable by any residual-based method regardless of
the detection threshold. Non-linear detectors based on learned
representations can potentially catch it if the perturbation falls
outside the data manifold, but an adaptive attacker will minimise
the detector output directly.

\subsection{Replay Attacks}

A \emph{replay attack} on sensor $i$ at time $t$ sets
$\tilde{p}_i(t) = p_i^*(t - \tau)$ for some lag $\tau \geq 1$. The
replayed value is a genuinely observed past reading and therefore
lies on the data manifold; a spatial detector will find it perfectly
plausible. Only a temporal detector that compares the current reading
against the recent history of the same sensor can identify that the
reading has stopped evolving.

This is why temporal modelling is not just an incremental improvement
over spatial modelling: it is qualitatively necessary for replay
detection. A model that processes one snapshot at a time cannot
distinguish a replay from a genuine steady-state reading.

\section{Graph Neural Networks}
\label{sec:gnn_background}

\subsection{Message Passing Framework}

Graph neural networks learn node representations by iteratively
aggregating messages from neighbours~\cite{gilmer2017neural}. At
layer $l$, node $i$ updates its embedding by:
\begin{align}
  \mathbf{m}_i^{(l)} &= \text{AGGREGATE}\bigl(\{(\mathbf{h}_j^{(l-1)},
    \mathbf{e}_{ij}) : j \in \mathcal{N}(i)\}\bigr), \label{eq:aggregate} \\
  \mathbf{h}_i^{(l)} &= \text{UPDATE}\bigl(\mathbf{h}_i^{(l-1)},
    \mathbf{m}_i^{(l)}\bigr). \label{eq:update}
\end{align}
Different choices of AGGREGATE and UPDATE define different \gnn{}
architectures. After $L$ layers, $\mathbf{h}_i^{(L)}$ captures
information from the $L$-hop neighbourhood of $i$.

\subsection{Architectures Compared in This Thesis}

This thesis compares four backbone architectures for the \wdn{}
anomaly detection task.

\textbf{Graph Convolutional Network (GCN).} Kipf and Welling's
Graph Convolutional Network (GCN)~\cite{kipf2017semi} uses a symmetric normalised
aggregation:
\begin{equation}
  \mathbf{H}^{(l+1)} = \sigma\!\left(\hat{\mathbf{D}}^{-\frac{1}{2}}
    \hat{\mathbf{A}}\hat{\mathbf{D}}^{-\frac{1}{2}}
    \mathbf{H}^{(l)}\mathbf{W}^{(l)}\right),
\end{equation}
where $\hat{\mathbf{A}} = \mathbf{A} + \mathbf{I}$ is the adjacency
matrix with self-loops and $\hat{\mathbf{D}}$ is the corresponding
degree matrix. The normalisation prevents exploding or vanishing
activations on irregular-degree graphs but constrains the expressive
power relative to architectures with learnable aggregation weights.

\textbf{Graph Attention Network (GAT).} Veli{\v{c}}kovi{\'c}
et al.~\cite{velickovic2018graph} replace fixed normalisation with
learned attention coefficients:
\begin{equation}
  \alpha_{ij} = \frac{\exp\!\left(\text{LeakyReLU}\!\left(
    \mathbf{a}^\top[\mathbf{W}\mathbf{h}_i \| \mathbf{W}\mathbf{h}_j]
    \right)\right)}{\sum_{k \in \mathcal{N}(i)}
    \exp\!\left(\text{LeakyReLU}\!\left(
    \mathbf{a}^\top[\mathbf{W}\mathbf{h}_i \| \mathbf{W}\mathbf{h}_k]
    \right)\right)}.
\end{equation}
Multi-head attention averages $H$ such transformations. GAT
learns which neighbours to attend to, which is useful when some pipes
carry much more flow than others.

\textbf{GraphSAGE.} Hamilton et al.~\cite{hamilton2017inductive}
aggregate neighbour embeddings with a learned weight and concatenate
the result with the node's own embedding:
\begin{equation}
  \mathbf{h}_i^{(l+1)} = \sigma\!\left(\mathbf{W}^{(l)}\!\left[
    \mathbf{h}_i^{(l)} \,\big\|\, \text{MEAN}_{j \in \mathcal{N}(i)}
    \mathbf{h}_j^{(l)}\right]\right).
\end{equation}
The concatenation preserves the node's own identity across layers, which
matters in a \wdn{} where each junction has a distinct physical role.
The GNN backbone ablation in Chapter~\ref{ch:results} shows
GraphSAGE as the best-performing backbone on all three networks,
and it is used as the default throughout this thesis.

\textbf{Graph Transformer.} The transformer-based \gnn{}~\cite{shi2021masked}
applies multi-head self-attention over the neighbourhood, treating
node pairs as query-key pairs with positional encodings. It has the
highest parameter count in our ablation (around 55\,000 parameters
versus 29\,000 for GCN) but does not outperform
GraphSAGE on the water network tasks, likely because the local
neighbourhood structure is more informative than global attention
patterns on these relatively small graphs.

\subsection{Edge Features in Water Networks}

Pipe attributes — length, diameter, roughness coefficient, and
connectivity type (pipe, pump, valve) — carry information relevant
to how pressures propagate. Most early applications of \gnn{}s to
water networks ignored edge features; the models in this thesis
include them explicitly. At each message-passing layer the edge
feature vector $\mathbf{e}_{ij}$ is concatenated with the source
node's embedding before the aggregation weight is computed.

\section{Temporal Sequence Modelling}
\label{sec:temporal_background}

\subsection{Gated Recurrent Units}

The Gated Recurrent Unit (GRU)~\cite{cho2014gru} processes a sequence of inputs
$(x_1, \ldots, x_T)$ through an update gate $z_t$ and a reset gate
$r_t$:
\begin{align}
  z_t &= \sigma(\mathbf{W}_z x_t + \mathbf{U}_z h_{t-1}), \\
  r_t &= \sigma(\mathbf{W}_r x_t + \mathbf{U}_r h_{t-1}), \\
  \tilde{h}_t &= \tanh(\mathbf{W} x_t + \mathbf{U}(r_t \odot h_{t-1})), \\
  h_t &= (1 - z_t) \odot h_{t-1} + z_t \odot \tilde{h}_t.
\end{align}
The update gate interpolates between the previous hidden state and
the candidate update, allowing the \gru{} to retain long-range
dependencies without the gradient problems of vanilla RNNs. In this
thesis, a single \gru{} layer operates on the sequence of per-node
\gnn{} embeddings across a sliding window of $T = 6$ steps.

\subsection{Temporal Stability Features}

Beyond the learned \gru{} representation, nine hand-crafted temporal
features are appended to the node feature vector at each time step.
These are motivated by the observation that different attack types
leave characteristic statistical signatures in the recent history:

\begin{itemize}
  \item \emph{Autocorrelation} at lags 1 and 2: a replay attack
    replaces a changing reading with a constant one, driving
    autocorrelation to 1.0.
  \item \emph{Differential standard deviation}: the spread of
    first differences; noise attacks inflate this significantly.
  \item \emph{Rolling mean and standard deviation} over the window:
    a stealthy bias drift shifts the mean monotonically.
  \item \emph{Noise ratio}: the fraction of the variance
    attributable to high-frequency components; useful for jamming
    attacks.
  \item \emph{Trend slope}: the slope of a linear fit to the window;
    a stealthy attack has a non-zero, consistent slope.
\end{itemize}

Section~\ref{sec:pattern_features} in Chapter~\ref{ch:architecture}
quantifies the lift these features provide for each attack type.

\section{Mixture of Experts}
\label{sec:moe_background}

A Mixture of Experts (MoE) model routes each input to one of $K$ specialist
sub-networks (\emph{experts}) via a learned \emph{router}~\cite{jacobs1991moe,
shazeer2017outrageously}. The final output is a weighted combination
of expert outputs:
\begin{equation}
  \hat{y} = \sum_{k=1}^K g_k(\mathbf{x}) \, f_k(\mathbf{x}),
\end{equation}
where $g_k(\mathbf{x}) = \mathrm{softmax}(\text{Router}(\mathbf{x}))_k$
are the routing probabilities and $f_k$ is the $k$-th expert.

A key training challenge is \emph{mode collapse}: the router tends
to converge on a single expert and stop exploring alternatives.
The standard remedy is a load-balancing loss that encourages the
router to spread its probability mass:
\begin{equation}
  \mathcal{L}_\text{balance} = -\sum_{k=1}^K \bar{g}_k \log \bar{g}_k,
\end{equation}
where $\bar{g}_k$ is the batch-mean routing probability for expert
$k$. This term is added to the total loss with a small weight
$\lambda_\text{balance}$.

For the \moe{} defender in this thesis, the router uses the same
temporal \gnn{} architecture as the experts but with fewer layers.
The routing target during training is the ground-truth attack type,
which provides a supervision signal beyond the load-balancing entropy.
After training, the router correctly identifies the attack type in
74\% of snapshots on Net3.

\section{Adversarial Machine Learning and Self-Play}
\label{sec:adversarial_background}

\subsection{Adversarial Examples}

Goodfellow et al.~\cite{goodfellow2015explaining} showed that adding
a small, structured perturbation to an input can cause a
well-trained classifier to output an incorrect label with high
confidence. The \emph{fast gradient sign method} constructs such
a perturbation as $\boldsymbol{\delta} = \varepsilon \cdot
\mathrm{sgn}(\nabla_{\mathbf{x}} \mathcal{L})$, where $\mathcal{L}$
is the model's loss on the correct label. The projection-based variant
(PGD, Madry et al.~\cite{madry2018pgd}) iterates this step within an
$\varepsilon$-ball.

These attacks assume white-box access to the model. In a \wdn{}
context the analogous assumption is that the attacker knows the
defender architecture and weights and can compute gradients through
it.

\subsection{Adversarial Training}

The most direct defence against adversarial examples is adversarial
training~\cite{madry2018pgd}: augment the training set with adversarial
examples generated by an attack procedure and retrain. The difficulty
is that the attack procedure is fixed, so the model learns to resist
the specific attack strategy used during training but may generalise
poorly to a different strategy.

\subsection{Self-Play and Stackelberg Games}

Self-play is the idea of using the model's own adversarial copies as
training opponents. The concept became widely known through AlphaGo
Zero~\cite{silver2017alphagozero}, which trained a board game agent
by playing against earlier versions of itself without using any
human data. Applied to security, self-play produces a defender that
has, by construction, been tested against the strongest attacker the
defender's current weights can generate.

A \emph{Stackelberg game} is a sequential game where one player
(the \emph{leader}) commits to a strategy first and the other
(the \emph{follower}) best-responds~\cite{von2010leadership}. In
the adversarial training context, the attacker acts as the leader
(it is updated to exploit the current defender) and the defender
acts as the follower (it is updated to catch the current attacker).
This alternating structure is exactly what the training loop in
Chapter~\ref{ch:selfplay} implements.

\subsection{Mixture of Experts as an Attacker}

Extending the attacker to a \moe{} of specialists is the key
technical contribution that distinguishes this thesis from standard
adversarial training. The motivation is population diversity: a
single attacker converges to a single attack strategy once the
defender can resist it. A population of attackers, trained with a
pairwise diversity loss, can maintain multiple attack strategies
simultaneously. The result is a defender trained against a richer
threat distribution, which explains the improved performance on
held-out attack types.

\section{Related Work}
\label{sec:related}

\subsection{Machine Learning for WDN Anomaly Detection}

Early work on anomaly detection in water networks used statistical
process control~\cite{klise2006statistical} or autoencoder-based
methods that ignored graph structure~\cite{taormina2019deep}.
Hajgato et al.~\cite{hajgato2021pressure} applied a purely spatial
\gnn{} to pressure field reconstruction but did not address anomaly
detection. Zanfei et al.~\cite{zanfei2022graph} used a GCN
for leak detection but modelled only a single type of fault. None of
these works considered adversarial adaptation.

The most closely related prior work is the BattLeDIM
competition~\cite{vrachimis2021battledim}, which benchmarked leak
detection algorithms on the L-Town network. The winning solution used
a combination of hydraulic simulation residuals and statistical
tests; no \gnn{}-based solution was submitted. The anomaly detection
problem in this thesis is harder because the attacker is adversarially
optimal rather than following a fixed rule.

\subsection{GNNs for Cyber-Physical System Security}

Li et al.~\cite{li2021gnn} applied message-passing \gnn{}s to anomaly
detection in industrial control systems but worked on small, toy
graphs. Mathur et al.~\cite{mathur2016swat} introduced the Secure
Water Treatment (SWaT) dataset, which has been used for evaluating
detection algorithms on a water treatment plant; the graph structure
of the plant was not exploited in most follow-up work.

\subsection{Adversarial Self-Play for Security}

Apruzzese et al.~\cite{apruzzese2023role} survey adversarial ML
for cyber-security and note that most papers evaluate on fixed
attack benchmarks rather than adaptive attackers. Ilyas
et al.~\cite{ilyas2019adversarial} show that adversarial robustness
and standard accuracy are often in tension, which is what we see on
Net1.

The pieces this thesis combines exist separately in the literature.
Stackelberg games have been used for water-network security, but
in the resource-allocation sense (where to place defences), not
with a learned attacker model. Adversarial attacks on GNNs are
well studied on generic graph data. Temporal GCNs have been applied
to attack detection on the C-Town BATADAL
benchmark~\cite{tsiami2021cyber}. What we did
not find is the specific combination used here: a learned
attacker GNN, made into a mixture of specialists, co-trained
against the defender in an alternating Stackelberg loop, on water
distribution sensor data. We make the novelty claim only for that
combination, and only as far as our literature search reached.
